\documentclass[11pt,oneside]{article}
\usepackage[T1]{fontenc}
\usepackage[utf8]{inputenc}
\usepackage{lmodern}
\usepackage[a4paper,margin=27mm,headheight=15pt]{geometry}
\usepackage{microtype}
\usepackage{amsmath,amssymb,amsthm,mathtools,aliascnt}
\usepackage{booktabs,longtable,array,enumitem}
\usepackage[dvipsnames]{xcolor}
\usepackage{fancyhdr}
\usepackage{hyperref}
\hypersetup{colorlinks=true,linkcolor=MidnightBlue,citecolor=MidnightBlue,
 urlcolor=MidnightBlue,pdftitle={Omnific Integers and Omnific--Diophantine Geometry},
 pdfauthor={OpenAI ChatGPT},pdfsubject={Arithmetic, rigidity, definability, and infinite solutions over the omnific integers}}
\usepackage[nameinlink,noabbrev]{cleveref}
\pagestyle{fancy}
\fancyhf{}
\fancyhead[L]{\small Omnific integers and Diophantine geometry}
\fancyhead[R]{\small\thepage}
\renewcommand{\headrulewidth}{0.3pt}
\setlength{\parskip}{4pt}
\setlength{\parindent}{1.2em}
\setlength{\emergencystretch}{2em}
\setcounter{tocdepth}{2}
\numberwithin{equation}{section}
\newtheorem{theorem}{Theorem}[section]
\newaliascnt{proposition}{theorem}
\newtheorem{proposition}[proposition]{Proposition}
\aliascntresetthe{proposition}
\crefname{proposition}{proposition}{propositions}
\newaliascnt{lemma}{theorem}
\newtheorem{lemma}[lemma]{Lemma}
\aliascntresetthe{lemma}
\crefname{lemma}{lemma}{lemmas}
\newaliascnt{corollary}{theorem}
\newtheorem{corollary}[corollary]{Corollary}
\aliascntresetthe{corollary}
\crefname{corollary}{corollary}{corollaries}
\theoremstyle{definition}
\newaliascnt{definition}{theorem}
\newtheorem{definition}[definition]{Definition}
\aliascntresetthe{definition}
\crefname{definition}{definition}{definitions}
\newaliascnt{example}{theorem}
\newtheorem{example}[example]{Example}
\aliascntresetthe{example}
\crefname{example}{example}{examples}
\theoremstyle{remark}
\newaliascnt{remark}{theorem}
\newtheorem{remark}[remark]{Remark}
\aliascntresetthe{remark}
\crefname{remark}{remark}{remarks}
\newcommand{\No}{\mathbf{No}}
\newcommand{\Oz}{\mathbf{Oz}}
\newcommand{\Z}{\mathbb Z}
\newcommand{\Q}{\mathbb Q}
\newcommand{\R}{\mathbb R}
\newcommand{\C}{\mathbb C}
\newcommand{\N}{\mathbb N}
\newcommand{\J}{\mathcal J}
\newcommand{\BR}{\mathcal B_{\R}}
\newcommand{\BC}{\mathcal B_{\C}}
\newcommand{\ct}{\operatorname{ct}}
\newcommand{\supp}{\operatorname{supp}}
\newcommand{\lc}{\operatorname{lc}}
\newcommand{\ddeg}{\operatorname{deg}_{\omega}}
\newcommand{\Frac}{\operatorname{Frac}}
\newcommand{\Int}{\operatorname{Int}}
\newcommand{\rad}{\operatorname{rad}}
\newcommand{\GL}{\operatorname{GL}}
\newcommand{\diag}{\operatorname{diag}}
\newcommand{\tr}{\operatorname{tr}}
\newcommand{\id}{\operatorname{id}}
\newcommand{\OO}{\operatorname O}
\newcommand{\V}{\operatorname V}
\newcommand{\repo}{\texttt{VladimirReshetnikov/Surreal}}
\newcommand{\pin}{\texttt{2cb9c0200af749fbbd27796c97d017e7f16fbf33}}
\newcommand{\status}[1]{\par\smallskip\noindent\textit{Provenance.} #1\par\smallskip}

\begin{document}
\begin{titlepage}
\centering
\vspace*{16mm}
{\Large Research article for the Surreal project\par}
\vspace{13mm}
{\huge\bfseries Omnific Integers and\\[3mm] Omnific--Diophantine Geometry\par}
\vspace{8mm}
{\Large Retractions, rigidity, definability,\\ and infinite families\par}
\vspace{14mm}
{\large September 22, 2026\par}
\vspace{10mm}
\begin{minipage}{0.91\textwidth}
\begin{abstract}
The omnific integers form a discretely ordered subring of the surreal field,
but their arithmetic is neither ordinary integer arithmetic nor a routine
nonstandard enlargement of it. We organize their theory around the split
ring homomorphism that takes a normal form to its constant coefficient.
This gives a precise transfer theorem for polynomial equations with ordinary
integer coefficients, a classification of finite quotients, and a complete
solution criterion for linear systems.

The main Diophantine results distinguish existence from the structure of
solutions. Nonzero fibers of nondegenerate binary forms and number-field
norm forms have no new omnific points. In particular, Pell equations admit
only ordinary-integer solutions. In contrast, every represented nonzero
level of an indefinite nondegenerate quadratic form in at least three
variables has explicit infinite polynomial families through each integer
point. We prove the resulting quadratic-level dichotomy.

Combining Pell rigidity with the classical three-square theorem produces
a single quartic polynomial, with five auxiliary variables independent of
the tuple length, defining exactly the ordinary integers inside the omnific
integers. Further results concern non-atomic factorization, a common-divisor
phenomenon peculiar to the full surreal class, projective denominator
clearing, primitive Pythagorean triples, homogeneous cones, and obstructions
at infinity. All central deductions receive proofs. Historical background,
imported theorems, proposed contributions, and unverified formalization work
are explicitly separated.
\end{abstract}
\end{minipage}
\vfill
\begin{minipage}{0.91\textwidth}
\small
Prepared with OpenAI ChatGPT. This is a mathematical research draft, not a
claim of peer review or proof-assistant verification. The included program
checks finite identities and examples, not the entire theory. Repository
inspection is pinned to the commit recorded in Appendix~\ref{app:audit}.
\end{minipage}
\end{titlepage}
\tableofcontents
\newpage

\section{Scope, results, and mathematical conventions}\label{sec:scope}

Omnific integers are the surreal numbers with no negative exponents in
Conway normal form and with an ordinary integer constant coefficient.
The definition is classical, as are their discreteness and integer-part
property \cite{Conway,Gonshor,LM}. This article develops their arithmetic
with a particular emphasis on polynomial equations.

Three questions must be kept separate. Does an equation have any solution?
Does it have a genuinely infinite solution? Does it have such a solution
with a primitivity condition? For example, every system of equations over
$\Z$ has an omnific solution exactly when it has an ordinary integer
solution. Nevertheless, the equation $X^n+Y^n=Z^n$ has positive infinite
omnific solutions for every ordinary $n\geq2$. There is no contradiction:
the constant coefficients of these solutions are all zero.

We write $A=\Oz$. The ideal $\J$ consists of the normal forms supported on
strictly positive exponents; coefficients in $\J$ may have either sign.
The central decomposition is
\begin{equation}\label{eq:splitintro}
 A=\Z\oplus\J,\qquad \ct:A\longrightarrow\Z,
 \qquad \ct(a+b)=\ct(a)+\ct(b),\quad \ct(ab)=\ct(a)\ct(b).
\end{equation}
It is a splitting of additive groups and a ring retraction, not a direct
product decomposition of rings. In particular, $\J$ has no identity.

\subsection{Principal results}

The structural results in Sections~\ref{sec:basic}--\ref{sec:class} explain
why ordinary congruences see only the constant coefficient, although the
fraction field of the full class $A$ is all of $\No$. The main Diophantine
statements are the following.

\begin{enumerate}[label=\textup{(\arabic*)},leftmargin=*,itemsep=3pt]
\item \textbf{Equational transfer.} For every finite system over $\Z$,
$\V(A)$ is nonempty exactly when $\V(\Z)$ is nonempty
(\cref{thm:transfer}). This does not include inequations or order
conditions.
\item \textbf{Constant-product rigidity.} A nonzero constant product of
normal forms with no negative exponents has only constant factors. This
forces all omnific points on nonzero binary-form and number-field norm
fibers to be ordinary integer points
(\cref{thm:binary,thm:norm}).
\item \textbf{Quadratic-level dichotomy.} For $q\in\Z[X_1,\ldots,X_n]$
quadratic homogeneous and $0\ne c\in\Z$, all solutions are ordinary in
the nondegenerate definite case and in nondegenerate dimension at most two.
In every other case, each integer point lifts to an injectively
parameterized infinite family (\cref{thm:quadclassification}).
\item \textbf{A quartic definition of $\Z^n$.} The equation
\[
 (U^2-2V^2-1)^2+
 \left(U-\sum_{i=1}^{n}X_i^2-W_1^2-W_2^2-W_3^2\right)^2=0
\]
with five existentially quantified auxiliary variables defines exactly
$\Z^n$ inside $A^n$ (\cref{thm:guard}).
\item \textbf{Primitive and projective phenomena.} There are infinite
unimodular Pythagorean triples. On the other hand, a constant real
projective direction has a primitive omnific representative precisely when
it is rational (\cref{prop:pythagoras,thm:realray}).
\end{enumerate}

These statements are proved here, but proving a statement here is not a
claim that nobody has previously proved it. The constant-term formalism is
classical; several consequences are natural applications of it. The
quadratic classification and explicit quartic guard are the most distinctive
constructions developed in this article. The inspected sources did not
provide these statements in this form. No exhaustive historical-priority
claim is made.

\subsection{Class size and finite algebra}

The full surreal field and $A$ are proper classes. We work in a customary
class-theoretic setting, such as NBG with the usual background for surreal
normal forms. Each individual surreal has a \emph{set-sized} support.
Every polynomial, system, matrix, or tuple in this article has ordinary
finite size. Quantifiers over $A$ are understood as class quantifiers
expressed through their defining set-theoretic formulas, not as an appeal
to a set containing all surreals.

Most proofs can instead be carried out in a fixed set-sized Hahn field.
We explicitly identify the arguments that need the stronger fact that
\emph{every set} of surreal exponents has upper bounds and every set of
positive surreal exponents has a strictly positive lower bound. One must
not silently assume these properties for a fixed ordered exponent group.

The notation $\omega^\gamma$ always means the Conway monomial. Its
multiplication is $\omega^\gamma\omega^\delta=\omega^{\gamma+\delta}$,
where $+$ is surreal addition, not ordinary ordinal addition. Exponent
supports are reverse well ordered. In the frequent Hahn convention
$t=\omega^{-1}$, all exponent signs are reversed.

\subsection{Relationship with the repository and literature}

The repository's documentation distinguishes normal-form algebra, fixed
Hahn workspaces, and the full surreal carrier. Its existing trigonometry
material already uses the omnific integer part and residue arithmetic
\cite{Repo}. We retain those distinctions and develop a separate
Diophantine treatment rather than identifying $A$ with a finite-surreal
valuation ring.

The main external sources are Conway and Gonshor for the construction and
normal forms \cite{Conway,Gonshor}, L'Innocente--Mantova for modern
factorization theory \cite{LM}, Ehrlich--Kaplan for the initial integer-part
context \cite{EK}, and Je\v r\'abek for the general residue theory of
$\Z$-groups \cite{Jerabek}. The three-square theorem and MRDP theorem are
imported with explicit attribution at their points of use.

\section{The normal-form ring and its integer-part structure}\label{sec:basic}

\subsection{Normal forms and three related rings}

We use the standard normal-form theorem: each surreal has a unique
expression
\begin{equation}\label{eq:nf}
 x=\sum_{\gamma\in S}r_\gamma\omega^\gamma,
 \qquad r_\gamma\in\R\setminus\{0\},
\end{equation}
where $S\subseteq\No$ is a reverse-well-ordered set. Conversely such a
normal form represents a surreal. Addition is coefficientwise,
multiplication is the well-defined Hahn convolution, and the sign of a
nonzero normal form is the sign of its leading coefficient
\cite{Conway,Gonshor}.

\begin{definition}\label{def:rings}
Let
\begin{align*}
 \J&=\{x\in\No:\supp(x)\subseteq\No^{>0}\},\\
 \BR&=\R\oplus\J,\\
 A=\Oz&=\Z\oplus\J,\\
 \BC&=\BR[i]\subseteq\No(i).
\end{align*}
The zero series belongs to $\J$. For an element of any of these rings,
$\ct(x)$ denotes its coefficient at exponent zero. For
$x=a+ib\in\BC$, this is $\ct(a)+i\ct(b)$.
\end{definition}

Thus $\sqrt2\,\omega$, $\omega/17$, $-\omega^\omega+3$, and
$\sum_{k\geq1}\omega^{1/k}$ are omnific integers. The numbers
$\omega+\tfrac12$ and $\omega+\omega^{-1}$ are not. The word ``integer''
restricts the constant coefficient and excludes infinitesimal tails; it
does not require all coefficients to lie in $\Z$.

For every positive surreal exponent $\gamma$, even an infinitesimal one,
the Conway monomial $\omega^\gamma$ exceeds every ordinary positive
real number. This is a normal-form order fact about the omega-map;
replacing that map by an unstated analytic notion of exponentiation would
invalidate the support arguments below.

\begin{proposition}\label{prop:ring}
$\BR$, $\BC$, and $A$ are integral domains. The class $\J$ is an ideal
of $\BR$ and $A$, and is an $\R$-vector space. Constant-coefficient
extraction is a surjective ring homomorphism from $\BR$ to $\R$, from
$\BC$ to $\C$, and from $A$ to $\Z$. In the last case its kernel is
$\J$, and the inclusion $\Z\hookrightarrow A$ is a section.
\end{proposition}
\begin{proof}
Nonnegative exponents remain nonnegative under addition of exponents. A
strictly positive exponent plus a nonnegative exponent is strictly
positive. These observations and Hahn convolution prove the ring and ideal
claims. In a product of two supported normal forms, exponent zero can
arise only as $0+0$. Therefore $\ct(xy)=\ct(x)\ct(y)$. The additive
claim is immediate. Surjectivity and the section are supplied by constant
series. All rings embed in the relevant field, so have no zero divisors.
\end{proof}

\begin{remark}[A frequently misleading analogy]\label{rem:notvaluation}
The map $\ct$ is not a ring homomorphism on all of $\No$:
$\ct(\omega)=\ct(\omega^{-1})=0$ but $\ct(1)=1$. Also $\BR$ is not
the ring of finite surreals. The latter permits negative exponents but no
positive exponents; $\BR$ has the opposite support restriction. In
particular, taking standard part of a finite element and taking constant
coefficient of an omnific integer are different constructions, although
both are multiplicative on their own respective domains.
\end{remark}

\subsection{Leading degree, units, and finite elements}

For nonzero $x\in\BC$, combine real and imaginary coefficients and put
\[
 \ddeg(x)=\max\supp(x),\qquad \lc(x)=[\omega^{\ddeg(x)}]x.
\]
The degree lies in $\No^{\geq0}$. It is not an ordinal and not a birthday.

\begin{lemma}\label{lem:degree}
For nonzero $x,y\in\BC$,
\[
 \ddeg(xy)=\ddeg(x)+\ddeg(y),\qquad
 \lc(xy)=\lc(x)\lc(y).
\]
An element of $\BC$ of degree zero is an ordinary complex constant.
\end{lemma}
\begin{proof}
If $\gamma$ and $\delta$ are the respective largest exponents, every
contributing exponent is at most $\gamma+\delta$. Equality requires both
exponents to be largest: making either smaller makes their sum smaller.
Thus the coefficient at $\gamma+\delta$ is the nonzero product of leading
coefficients. Finally, support in $[0,0]$ contains only zero.
\end{proof}

\begin{proposition}\label{prop:units}
The unit groups are
\[
 \BR^\times=\R^\times,\qquad
 \BC^\times=\C^\times,\qquad A^\times=\{1,-1\}.
\]
Moreover, the omnific integers bounded in absolute value by some ordinary
real number are exactly the ordinary integers. Consequently $A$ is
discretely ordered, with least positive element $1$.
\end{proposition}
\begin{proof}
If $xy=1$ in $\BC$, nonnegativity and additivity of degree force both
degrees to be zero. This gives the first two unit statements; in $A$ the
constant units are precisely the units of $\Z$.

A nonconstant element of $A$ has positive leading exponent, so its
absolute value exceeds every positive real number by normal-form order.
A constant element of $A$ is an integer. Hence no element lies strictly
between $0$ and $1$, and every positive element is at least $1$.
\end{proof}

Notice that $\omega/2$ is an omnific integer even though $1/2$ is not.
This is no conflict with discreteness: $\omega/2$ is infinite, and
there are still no omnific integers in the interval $(0,1)$.

\subsection{The exact omnific floor}

Every surreal admits a unique decomposition
\[
 x=P+r+\varepsilon,
 \qquad P\in\J,\quad r\in\R,\quad
 \supp(\varepsilon)\subseteq\No^{<0}.
\]
Here $\varepsilon$ is infinitesimal, possibly zero.

\begin{theorem}[Integer-part formula]\label{thm:floor}
There is exactly one $a\in A$ such that $a\leq x<a+1$. It is
\begin{equation}\label{eq:floor}
 \lfloor x\rfloor_A=
 \begin{cases}
 P+\lfloor r\rfloor,&r\notin\Z\text{ or }\varepsilon\geq0,\\
 P+r-1,&r\in\Z\text{ and }\varepsilon<0.
 \end{cases}
\end{equation}
\end{theorem}
\begin{proof}
For $r\notin\Z$, its distances from the two adjacent integers are
strictly positive real numbers. The infinitesimal $\varepsilon$ cannot
cross either boundary. If $r$ is an integer, the sign of $\varepsilon$
determines on which side of $r$ the number lies, and its magnitude is
less than $1$. This proves the displayed inequalities. Two possible
floors would differ by an omnific integer strictly between $-1$ and $1$,
so would coincide by discreteness.
\end{proof}

For example,
\[
 \lfloor\omega-\omega^{-1}\rfloor_A=\omega-1,
 \qquad \lfloor\omega+\tfrac13-\omega^{-1}\rfloor_A=\omega.
\]
Merely discarding negative exponents gives the wrong answer at a negative
infinitesimal displacement from an integer.

\status{The integer-part assertion is classical. Ehrlich--Kaplan prove the
stronger initiality characterization: for an initial subfield $K$ of
$\No$, $\Oz\cap K$ is its unique initial integer part
\cite[Lemma 11.1]{EK}. No alternative integer part is being introduced here.}

\section{Divisibility, residues, and finite quotients}\label{sec:residues}

\subsection{Division by ordinary integers}

\begin{theorem}\label{thm:finitequotients}
For every nonzero ordinary integer $n$,
\begin{equation}\label{eq:nA}
 nA=\J+n\Z,\qquad
 n\mid a\ \Longleftrightarrow\ n\mid\ct(a),\qquad
 A/nA\cong\Z/n\Z.
\end{equation}
In particular,
\begin{equation}\label{eq:divkernel}
 \bigcap_{n\geq1}nA=\J,
 \qquad \bigcap_{k\geq1}p^kA=\J
 \quad\text{for every ordinary prime }p.
\end{equation}
\end{theorem}
\begin{proof}
Multiplication by a nonzero real scalar preserves $\J$ and is invertible
on $\J$. Thus $n\J=\J$. If $a=m+j$, then $a/n=m/n+j/n$ lies in
$A$ exactly when $m/n$ is an integer. Reduction of $m$ modulo $n$ gives
the quotient isomorphism. Finally, the only ordinary integer divisible by
all positive integers, or by every power of one prime, is zero.
\end{proof}

This theorem solves every ordinary congruence. For instance,
$\sqrt2\,\omega+14$ is even and is congruent to $2$ modulo $3$.
The whole infinite part is invisible to these tests.

The statement does \emph{not} extend to arbitrary omnific divisors.
For example, $\ct(\omega+2)$ divides $\ct(2)$ in $\Z$, but $\omega+2$
does not divide $2$ in $A$. A product equal to the nonzero constant $2$
can have only constant factors by \cref{lem:degree}.

\begin{corollary}\label{cor:charideals}
Every proper ideal $I\subset A$ whose quotient has positive characteristic
is $nA$ for a unique ordinary $n\geq2$. Such an ideal is prime exactly
when $n$ is an ordinary prime. Every nonzero finite quotient of $A$ is
therefore a ring $\Z/n\Z$.
\end{corollary}
\begin{proof}
If $m\cdot1\in I$ with $m\ne0$, then $\J=m\J\subseteq I$.
Thus $I$ is the inverse image of an ideal of $A/\J=\Z$.
Its image is $n\Z$ for some positive $n$, and properness gives $n\geq2$.
Equation~\eqref{eq:nA} identifies its inverse image with $nA$.
Primality is equivalent to primality of $n\Z$. A finite nonzero unital
ring has positive characteristic, so the last assertion follows.
\end{proof}

More generally, every unital homomorphism from $A$ to a finite ring kills
$\J$: if the target has characteristic $m$, write $j=m(j/m)$. Its image
is therefore the cyclic subring generated by $1$. This describes maps
into finite rings as well as finite quotients.

\subsection{Completions and canonicality}

All finite-index ideals have just been classified, so the profinite
completion, interpreted as the set-sized inverse limit of these quotients,
is
\begin{equation}\label{eq:profinite}
 \widehat A\cong\varprojlim_n\Z/n\Z=\widehat\Z.
\end{equation}
The completion map is the composite $A\xrightarrow{\ct}\Z\to\widehat\Z$,
with kernel $\J$. Similarly,
\[
 \varprojlim_k A/p^kA\cong\Z_p
\]
has kernel $\J$. Thus these congruence topologies on $A$ are not
Hausdorff. The corresponding separated quotient is $A/\J=\Z$;
completion does not preserve the infinite part.

\begin{proposition}\label{prop:canonicalct}
The constant-term map is the unique unital ring homomorphism $A\to\Z$.
Every unital endomorphism $f:A\to A$ preserves $\J$ and satisfies
$\ct\circ f=\ct$.
\end{proposition}
\begin{proof}
For a homomorphism $h:A\to\Z$, the integer $h(j)$ is divisible by every
positive integer whenever $j\in\J$. Hence $h(j)=0$. Since $h$ fixes
ordinary integers, $h=\ct$. An endomorphism preserves divisibility by
each ordinary integer, so preserves their common divisibility kernel
\eqref{eq:divkernel}. It fixes $\Z$ pointwise, which proves the final
identity.
\end{proof}

\status{The divisible-kernel and residue-map perspective belongs to
classical $\Z$-group theory; see \cite[Sections 2 and 4]{Jerabek}.
Here it has an especially concrete ring-theoretic realization because
$\ct$ is multiplicative. The finite-quotient and completion statements
above are proved directly rather than imported as additional assumptions.}

\section{Global algebra: common divisors and the full fraction field}\label{sec:class}

The next results use all surreal exponents, not merely the exponents in one
fixed Archimedean or finite-rank group. This is an important source of
arithmetic differences between full omnific integers and restricted series
rings.

\subsection{One monomial divides a whole set}

\begin{lemma}[Set-sized exponent bounds]\label{lem:setbounds}
If $E\subset\No^{>0}$ is a set, there is a surreal $\epsilon>0$ such
that $\epsilon<e$ for every $e\in E$. Every set of surreals also has a
strict upper bound.
\end{lemma}
\begin{proof}
A surreal separating the small cut $\{0\mid E\}$ gives the first
assertion. A separator of $\{E\mid\varnothing\}$ gives the second.
The assertions also hold for the empty set. These are applications of the
standard set-cut property, not an assertion of an infimum or supremum.
\end{proof}

\begin{theorem}[Common monomial divisibility]\label{thm:commondivisor}
For every set $S\subseteq\J$ there is $\epsilon>0$ such that
\[
 \omega^\epsilon\mid s\quad\text{and}\quad
 s/\omega^\epsilon\in\J\qquad(s\in S).
\]
In particular, each nonzero element of $\J$ is a product of two nonunits
of $A$.
\end{theorem}
\begin{proof}
The union of the supports of the elements of $S$ is a set $E$ of strictly
positive surreal exponents. Choose $\epsilon$ as in
\cref{lem:setbounds}. Dividing by $\omega^\epsilon$ shifts every exponent
$e$ to $e-\epsilon>0$, preserving the support order and its set size.
The quotient is therefore in $\J$. For a nonzero $s$, both factors are
nonzero elements with positive degree, and neither is a unit.
\end{proof}

Even the series $\sum_{k\geq1}\omega^{1/k}$ has such a divisor. The
positive number $\epsilon$ may be infinitesimal relative to every real
$1/k$. This explains why the proof cannot just be copied into a Hahn
ring with exponent group $\R$.

\begin{corollary}\label{cor:Jglobal}
The ideal $\J$ is idempotent, $\J^2=\J$, and it is not generated as an
ideal by any set of elements. Also, no nonzero element of $\J$ is a
finite product of irreducibles of $A$.
\end{corollary}
\begin{proof}
The inclusion $\J^2\subseteq\J$ follows from the ideal property. For
$j\ne0$, the factorization in \cref{thm:commondivisor} has both factors
in $\J$, proving the reverse inclusion.

Suppose a set $S\subseteq\J$ generated $\J$. Choose $\epsilon>0$
strictly below every exponent in its supports. Each finite $A$-linear
combination of elements of $S$ has support strictly above $\epsilon$:
multiplication by an element of $A$ only adds nonnegative exponents.
But $\omega^{\epsilon/2}\in\J$ cannot have such a support. This is a
contradiction.

Finally, if $j=p_1\cdots p_k\ne0$ and every $p_i$ were irreducible,
then $0=\ct(j)=\prod_i\ct(p_i)$ would put some $p_i$ in $\J$.
That factor would be reducible by \cref{thm:commondivisor}.
\end{proof}

This is an explicit failure of atomic factorization, not merely a
failure of uniqueness of factorization. It leaves room for nonconstant
prime elements having nonzero constant coefficient. For example,
L'Innocente--Mantova prove that
$\omega^{\sqrt2}+\omega+1$ is prime \cite[Theorem B]{LM}. That theorem
is not needed in any subsequent proof here.

\begin{proposition}[A necessary residue condition for irreducibility]\label{prop:irreduciblect}
The constant irreducibles of $A$ are exactly the ordinary signed primes.
Every nonconstant irreducible has constant coefficient $1$ or $-1$.
The latter condition is not sufficient for irreducibility.
\end{proposition}
\begin{proof}
A factorization of a nonzero ordinary integer has only constant factors
by degree additivity, so ordinary primality decides constant
irreducibility. If a nonconstant element $a$ has $|\ct(a)|>1$, choose
an ordinary prime $p$ dividing its constant coefficient. Then
$a=p(a/p)$ is a factorization into two nonunits: the second factor
remains nonconstant. If $\ct(a)=0$, the factorization $a=2(a/2)$
has the same property. This leaves only constant coefficients $\pm1$.
Finally, $(\omega+1)^2$ has constant coefficient $1$ and is reducible.
\end{proof}

The factorization $j=2(j/2)$ already proves that a nonzero $j\in\J$
is reducible, even in a fixed exponent workspace. The stronger content of
\cref{thm:commondivisor} is the existence of a common \emph{infinite
monomial} divisor for an entire set, with all quotients still in $\J$.

\subsection{Clearing arbitrary set-sized families of denominators}

\begin{theorem}[Global monomial clearing]\label{thm:fractions}
For every set $S\subseteq\No$ there is $h>0$ such that
\[
 \omega^h\in\J,\qquad \omega^h x\in\J\quad(x\in S).
\]
Consequently $\Frac(A)=\No$, where the equality means that every surreal
is a quotient of two omnific integers with nonzero denominator.
\end{theorem}
\begin{proof}
Take the set $E$ obtained by uniting the supports of all elements of $S$.
Choose $h>0$ strictly greater than every $-e$ for $e\in E$.
Multiplication by $\omega^h$ shifts every exponent to $h+e>0$.
All shifted normal forms belong to $\J$. Applying this to $S=\{x\}$
gives $x=(\omega^h x)/\omega^h$. The reverse inclusion in the fraction
field identity holds because $A$ is a subring of $\No$.
\end{proof}

For example, a support unbounded below in a chosen real exponent group is
still a set of surreal exponents and can be shifted above zero by a new,
larger surreal exponent. The theorem does not assert that the required
monomial lies in the original fixed Hahn workspace.

\begin{corollary}\label{cor:projectiveclear}
Every point of $\mathbb P^n(\No)$ has homogeneous coordinates in
$\J^{n+1}$. More generally, any fixed set-sized collection of coordinate
tuples can be cleared by a single positive monomial.
\end{corollary}
\begin{proof}
Apply \cref{thm:fractions} to the set of all coordinates and multiply
all tuples by the resulting nonzero scalar.
\end{proof}

This is projective \emph{representability}, not primitivity. In fact the
coordinates in this construction always have a common nonunit divisor.
Section~\ref{sec:primitive} will show that some real projective directions
have no primitive omnific representative at all.

\subsection{Division with remainder is not a terminating Euclidean algorithm}

\begin{proposition}\label{prop:division}
For $a,b\in A$ with $b>0$, there are unique $q,r\in A$ such that
\[
 a=bq+r,\qquad 0\leq r<b.
\]
\end{proposition}
\begin{proof}
Take $q=\lfloor a/b\rfloor_A$ from \cref{thm:floor} and
$r=a-bq$. The floor inequalities give the required remainder bounds.
Conversely those bounds force $q\leq a/b<q+1$, so floor uniqueness
determines $q$ and then $r$.
\end{proof}

Thus the familiar division step exists. Termination of repeated division
is the missing ingredient: the positive order on $A$ is not well founded.
For an explicit example, start with
\[
 r_0=\sqrt2\,\omega,\qquad r_1=\omega,
 \qquad \alpha=\sqrt2-1\in(0,1).
\]
The first quotient is $1$ and the first remainder is
$r_2=\alpha\omega$. Every subsequent quotient is $2$, and
\[
 r_k=\alpha^{k-1}\omega\qquad(k\geq1).
\]
Indeed, $1/\alpha=\sqrt2+1$ has ordinary floor $2$, and
$1-2\alpha=\alpha^2$. Every displayed remainder is positive and
infinite, so the ordinary finite Euclidean procedure never terminates.

\begin{proposition}[An explicit pair without a gcd]\label{prop:nogcd}
The elements $\sqrt2\,\omega$ and $\omega$ have no greatest common
divisor in $A$, in the divisibility sense. In particular, $A$ is not a
gcd domain or a B\'ezout domain.
\end{proposition}
\begin{proof}
Let $d\ne0$ be a common divisor and put $\delta=\ddeg(d)$.
Degree additivity gives $0\leq\delta\leq1$. If $\delta=1$, both
quotients $\sqrt2\,\omega/d$ and $\omega/d$ are degree-zero
omnific integers, hence nonzero ordinary integers. Their ratio would make
$\sqrt2$ rational. Therefore $\delta<1$.

Put $\eta=(\delta+1)/2$, so $\delta<\eta<1$. The element
$\omega^\eta$ is also a common divisor, since both quotients have
strictly positive exponent $1-\eta$. But $\omega^\eta$ cannot divide
$d$ in $A$, as its degree is larger than $\delta$. Thus no proposed
common divisor is divisible by every common divisor, as a gcd would be.
A B\'ezout domain would supply a gcd for each pair, so that possibility
is excluded as well.
\end{proof}

\subsection{Non-Noetherian and non-normal behavior}

The strictly increasing chain
\[
 (\omega)\subsetneq(\omega^{1/2})\subsetneq
 (\omega^{1/4})\subsetneq\cdots
\]
already witnesses a failure of the ascending chain condition on ideals.
The inclusions hold because the quotients of consecutive generators have
positive exponents. They are strict because the reverse quotients have
negative exponents and are not in $A$.

\begin{proposition}\label{prop:notnormal}
The full ring $A$ is not integrally closed in its fraction field.
\end{proposition}
\begin{proof}
Since $\No$ is real closed, let $r=\sqrt{\omega^2+1}>0$.
By \cref{thm:fractions}, $r\in\Frac(A)$, and it is integral over $A$
because it satisfies the monic equation $T^2-(\omega^2+1)=0$.
If $r\in A$, then
\[
 (r-\omega)(r+\omega)=1.
\]
Both factors would be units of $A$, hence would lie in $\{1,-1\}$.
Their difference $2\omega$ would then be an ordinary integer, a
contradiction. Thus $r\notin A$.
\end{proof}

In particular, the full omnific ring is not a UFD or a PID. Its
factorization theory cannot be replaced by an ordinary Euclidean algorithm
or by an unstated gcd assumption. The modern theory in \cite{LM} is
substantially subtler. The refinement property is different from
integral closedness or atomicity; none of the preceding arguments refutes
it. Appendix~\ref{app:audit} records why we make no current open-status
claim concerning Conway's refinement problem.

\section{Equational transfer and linear Diophantine equations}\label{sec:transfer}

\subsection{Existence is exactly ordinary existence}

For a finite family of polynomials $F_1,\ldots,F_r$ over a ring $R$, write
$\V_R(F_1,\ldots,F_r)$ for its simultaneous zero set in $R^n$.

\begin{theorem}[Retraction theorem for equations]\label{thm:transfer}
Let $F_1,\ldots,F_r\in\Z[X_1,\ldots,X_n]$. Coordinatewise constant
coefficient gives a retraction
\[
 \V_A(F_1,\ldots,F_r)\longrightarrow
 \V_\Z(F_1,\ldots,F_r).
\]
In particular,
\begin{equation}\label{eq:existencetransfer}
 \V_A(F_1,\ldots,F_r)\ne\varnothing
 \quad\Longleftrightarrow\quad
 \V_\Z(F_1,\ldots,F_r)\ne\varnothing.
\end{equation}
The same equivalence holds for positive existential formulas in the ring
language, evaluated at ordinary integer parameters.
\end{theorem}
\begin{proof}
For every polynomial $F$ over $\Z$ and $\mathbf x\in A^n$,
\[
 \ct(F(\mathbf x))=F(\ct(x_1),\ldots,\ct(x_n)),
\]
because $\ct$ is a ring homomorphism fixing all coefficients. Thus the
constant tuple of an omnific solution is an integer solution. Conversely,
an integer solution is already an omnific solution, and applying $\ct$
to it does nothing.

A positive existential formula is built from equalities using finite
conjunction, finite disjunction, and existential quantification. Ring
homomorphisms preserve such formulas, as follows directly by induction on
their formation; witnesses are mapped coordinatewise. Apply this first to
$\ct$, and then to the inclusion of $\Z$ in $A$.
\end{proof}

The theorem answers a question about existence, not cardinalities or
complete solution sets. For $Y=X^2$, the fiber of the retraction above
$(0,0)$ already contains $(t,t^2)$ for every $t\in\J$.

\begin{corollary}[Hilbert's tenth problem]\label{cor:H10}
There is no algorithm deciding, from an ordinary integer polynomial,
whether it has a zero in $A$. The yes-instances are exactly the
ordinary-integer yes-instances of Hilbert's tenth problem, and form a
computably enumerable complete set under the usual computable reductions.
\end{corollary}
\begin{proof}
Equation~\eqref{eq:existencetransfer} identifies the two decision problems
on exactly the same finite inputs. The undecidability and completeness
of the ordinary problem are the Davis--Putnam--Robinson--Matiyasevich
(MRDP) theorem and its standard consequences \cite{Matiyasevich}.
Enumerating ordinary integer tuples semidecides the omnific yes-instances
as well; no enumeration of the proper class $A$ is involved.
\end{proof}

This corollary concerns \emph{ordinary integer coefficients}. A decision
problem allowing arbitrary omnific coefficients requires an explicit
finite representation scheme, and is a different question.

\subsection{Why signs and nonvanishing do not transfer}

Consider
\[
 (X+1)^2-2(Y+1)^2=0.
\]
It has no solution in $\N^2$, since $X+1$ and $Y+1$ would be positive
integers with rational ratio $\sqrt2$. However,
\[
 X=\sqrt2\,\omega-1,\qquad Y=\omega-1
\]
is a solution with both coordinates positive in $A$. Its constant tuple
is $(-1,-1)$, an ordinary solution of the equation without positivity
constraints.

\begin{proposition}\label{prop:notdiophantine}
With ordinary integer parameters, none of the subsets
\[
 A\setminus\{0\},\qquad A^{>0},\qquad A^{\geq0}
\]
is definable by a positive existential ring formula over $A$.
\end{proposition}
\begin{proof}
Every such definable subset must be stable under $\ct$ by the
homomorphism argument in \cref{thm:transfer}. But $\omega$ is nonzero
and positive while $\ct(\omega)=0$. Also $\omega-1>0$ whereas
$\ct(\omega-1)=-1<0$.
\end{proof}

Thus ``nonzero'' cannot be harmlessly added to an equational transfer
statement. Encoding nonvanishing by $XY=1$ would restrict $X$ to
$\{1,-1\}$, not to all nonzero omnific integers.

\subsection{A complete solution criterion for linear systems}

\begin{theorem}[Smith normal form over the omnific integers]\label{thm:smith}
Let $M\in\Z^{m\times n}$ and $\mathbf b\in A^m$. Choose unimodular
integer matrices $U,V$ with
\[
 UMV=D=\diag(d_1,\ldots,d_r,0),\qquad d_i>0,
\]
in the rectangular Smith normal form, and put $\mathbf c=U\mathbf b$.
Then $M\mathbf x=\mathbf b$ is solvable in $A^n$ if and only if
\begin{equation}\label{eq:smithcriterion}
 d_i\mid\ct(c_i)\quad(1\leq i\leq r),
 \qquad c_i=0\quad(r<i\leq m).
\end{equation}
All solutions are $\mathbf x=V\mathbf y$, where
\[
 y_i=c_i/d_i\quad(1\leq i\leq r),
 \qquad y_{r+1},\ldots,y_n\in A\text{ are arbitrary}.
\]
\end{theorem}
\begin{proof}
Unimodular integer matrices and their inverses act on omnific tuples.
The substitution $\mathbf x=V\mathbf y$ transforms the system into
$D\mathbf y=\mathbf c$. For each nonzero diagonal entry, division is
possible precisely under \cref{thm:finitequotients}. Zero rows require
the full equation $c_i=0$, not merely its constant coefficient. The
remaining coordinates are unconstrained.
\end{proof}

For ordinary $\mathbf b$, a useful equivalent description is
\[
 \mathbf x=\mathbf a+\mathbf u,
 \qquad \mathbf a\in\Z^n,\quad M\mathbf a=\mathbf b,
 \quad \mathbf u\in\J^n,\quad M\mathbf u=0.
\]
If $M$ has full column rank, all solutions are ordinary. For a concrete
one-equation example, every solution to $2x+3y=1$ is
\[
 x=-1+3t,\qquad y=1-2t,\qquad t\in A.
\]
The same formulas as over $\Z$ work, but the parameter class is larger.

\section{Rigidity of constant-product and norm equations}\label{sec:rigidity}

\subsection{The constant-product principle}

\begin{lemma}\label{lem:constantproduct}
If $b_1,\ldots,b_s\in\BC$ and
$b_1\cdots b_s=c\in\C^\times$, then every $b_i$ belongs to
$\C^\times$.
\end{lemma}
\begin{proof}
Each $b_i$ has inverse $c^{-1}\prod_{j\ne i}b_j$ in $\BC$.
Apply \cref{prop:units}. Equivalently, their nonnegative degrees have
sum zero, so every degree is zero.
\end{proof}

The importance of $\BC$ is that linear factors may have irrational or
complex coefficients. They need not themselves be omnific integers.
They do, however, have no negative exponents, so the unit argument applies.

\begin{theorem}[Decomposable polynomial fibers]\label{thm:decomposable}
Let
\[
 F(\mathbf X)=\lambda\prod_{j=1}^{s}L_j(\mathbf X),
 \qquad \lambda\in\C^\times,
\]
where the $L_j$ are affine-linear polynomials over $\C$. Suppose the
matrix of their linear parts has column rank $n$. If
$\mathbf x\in\BC^n$ and $F(\mathbf x)=c\in\C^\times$, then
$\mathbf x\in\C^n$. In particular, any such point in $A^n$ lies in
$\Z^n$.
\end{theorem}
\begin{proof}
By \cref{lem:constantproduct}, every $L_j(\mathbf x)$ is constant.
Select $n$ rows whose linear coefficient matrix is invertible over
$\C$. Subtract the constant terms and solve this ordinary linear system.
All coordinates of $\mathbf x$ are complex constants. Its intersection
with $A^n$ is $\Z^n$.
\end{proof}

The rank hypothesis is essential. The equation $X^d=1$ in two variables
forces $X$ to be constant but says nothing about $Y$.

\subsection{Binary forms, Pell equations, and power differences}

\begin{theorem}[Binary-form rigidity]\label{thm:binary}
Let $F\in\Z[X,Y]$ be a nonzero homogeneous binary form having at least
two distinct projective linear factors over $\C$. For each
$c\in\Z\setminus\{0\}$,
\[
 \{(x,y)\in A^2:F(x,y)=c\}
 =\{(x,y)\in\Z^2:F(x,y)=c\}.
\]
\end{theorem}
\begin{proof}
Factor $F$ into linear forms over $\C$, keeping their multiplicities.
Two distinct projective factors have linearly independent coefficient
vectors. The rank in \cref{thm:decomposable} is therefore two.
\end{proof}

This theorem includes the nonzero fibers of every irreducible binary form
of degree at least two over $\Q$. For irreducible forms of degree at least
three these are the usual Thue equations; the conclusion here is about
the absence of additional omnific solutions, not a new proof of ordinary
finiteness or an algorithm for the ordinary solutions.

\begin{corollary}[Pell rigidity]\label{cor:pell}
For ordinary integers $D\ne0$ and $c\ne0$, every solution in $A^2$ of
\[
 x^2-Dy^2=c
\]
is an ordinary integer solution. In particular,
\[
 \{(u,v)\in A^2:u^2-2v^2=1\}
 =\{(u,v)\in\Z^2:u^2-2v^2=1\}.
\]
\end{corollary}
\begin{proof}
Over $\C$, the two factors $X-\sqrt D\,Y$ and $X+\sqrt D\,Y$
are distinct because $D\ne0$. Apply \cref{thm:binary}.
\end{proof}

The classical Pell equation has ordinary solutions of arbitrarily large
finite size. It nevertheless has no single infinite omnific solution.
This distinction is a first warning against treating $A$ as an elementary
extension of ordinary integer arithmetic.

In the full surreal field, the contrast is visible without solving a
recurrence. For $D>0$, $c\ne0$, and $s\in\No^\times$,
\[
 x=\frac{s+c/s}{2},\qquad
 y=\frac{s-c/s}{2\sqrt D}
\]
satisfy $x^2-Dy^2=c$. Taking $s=\omega$ produces negative-exponent
terms, so neither coordinate is an omnific integer. Discarding those
terms changes the norm from $c$ to zero. A surreal solution cannot be
turned into an omnific solution by unverified truncation.

\begin{proposition}[The zero Pell fiber]\label{prop:zeropell}
If $D$ is a positive nonsquare ordinary integer, all solutions of
$x^2-Dy^2=0$ in $A^2$ are exactly
\[
 (x,y)=(\sqrt D\,t,t)\quad\text{or}\quad(-\sqrt D\,t,t),
 \qquad t\in\J.
\]
\end{proposition}
\begin{proof}
Factor in $\No$ to obtain $x=\pm\sqrt D\,y$. Write
$y=m+j$ with $m\in\Z$ and $j\in\J$. The constant coefficient of
$x$ is $\pm\sqrt D\,m$, which is an integer only if $m=0$.
Conversely, multiplication of $\J$ by a real scalar preserves $\J$,
so all displayed pairs belong to $A^2$ and solve the equation.
\end{proof}

Thus changing the right side from a nonzero constant to zero changes the
answer completely.

\begin{corollary}[Power-difference rigidity]\label{cor:powers}
Let $m,n\geq2$ be ordinary integers with $d=\gcd(m,n)>1$, and let
$0\ne c\in\Z$. Every omnific solution of $x^m-y^n=c$ is ordinary.
\end{corollary}
\begin{proof}
Set $X=x^{m/d}$ and $Y=y^{n/d}$. The binary form $X^d-Y^d$ has at
least two distinct complex factors, so $X,Y$ are ordinary by
\cref{thm:binary}. If a positive ordinary power of an element of $A$
is an ordinary constant, that element is ordinary: for a nonzero element
this follows from degree additivity, and for a zero power value the
integral domain property suffices.
\end{proof}

No conclusion for coprime exponents is being smuggled into this statement.
The argument does not classify equations such as $x^2-y^3=1$.

\subsection{Number-field norm equations}

\begin{theorem}[Norm-form rigidity]\label{thm:norm}
Let $K/\Q$ be a number field of degree $d$, let $e_1,\ldots,e_d$ be a
$\Q$-basis, and let
\[
 N(\mathbf X)=N_{K/\Q}\left(\sum_{j=1}^d e_jX_j\right).
\]
For every nonzero rational constant $c$, every solution
$\mathbf x\in A^d$ of $N(\mathbf x)=c$ has all coordinates in
$\Z$. In particular, with an integral basis this applies to the usual
integer-coefficient norm equations.
\end{theorem}
\begin{proof}
The $d$ embeddings $\sigma_1,\ldots,\sigma_d:K\to\C$ give
\[
 N(\mathbf X)=\prod_{i=1}^{d}
 \left(\sum_{j=1}^{d}\sigma_i(e_j)X_j\right).
\]
The embedding matrix $(\sigma_i(e_j))$ is invertible. To verify this,
choose a primitive element $\theta$ of the separable extension $K/\Q$.
For the power basis $1,\theta,\ldots,\theta^{d-1}$ the matrix is
Vandermonde in the distinct conjugates of $\theta$, so has nonzero
determinant. A basis change over $\Q$ preserves invertibility.
The result now follows from \cref{thm:decomposable}.
\end{proof}

For example, consider the norm form for $K=\Q(\sqrt[3]{2})$:
\[
 N(X+Y\sqrt[3]{2}+Z\sqrt[3]{4})
 =X^3+2Y^3+4Z^3-6XYZ.
\]
Every omnific solution of
\[
 X^3+2Y^3+4Z^3-6XYZ=1
\]
is an ordinary integer triple. This is a genuinely three-variable
rigidity result, unlike the indefinite quadratic forms in the next
section.

\subsection{An algebraic-geometric formulation}

Let $V$ be an affine complex variety whose coordinate ring $\C[V]$ is
generated, as a $\C$-algebra, by units. Every $\BC$-point of $V$ is
then a constant $\C$-point. Indeed, a $\C$-algebra map
$\C[V]\to\BC$ sends each generating unit to an element of
$\BC^\times=\C^\times$, so its entire image lies in $\C$.
This includes closed subvarieties of an algebraic torus, in their
very-affine coordinate presentation.

The statement supplies a conceptual explanation of the norm argument:
a nonzero product relation forces its factors into a constant torus.
It is not a claim that every affine variety is constant over $\BC$.
The affine line, for example, has the full class $\BC$ of points.

\section{A complete dichotomy for nonzero quadratic levels}\label{sec:quadratic}

Binary Pell rigidity might suggest that a nonzero constant quadratic level
never contains infinite points. This is false as soon as there is enough
indefinite dimension. The following result gives a precise classification.

\begin{theorem}[Quadratic-level dichotomy]\label{thm:quadclassification}
Let $n\geq1$, let $q\in\Z[X_1,\ldots,X_n]$ be a homogeneous quadratic
form, and let $c\in\Z\setminus\{0\}$. Write
\[
 V_A=\{\mathbf x\in A^n:q(\mathbf x)=c\},\qquad
 V_\Z=\{\mathbf a\in\Z^n:q(\mathbf a)=c\}.
\]
If $V_\Z$ is empty, so is $V_A$. Otherwise the following alternatives
are exhaustive.
\begin{enumerate}[label=\textup{(\roman*)},leftmargin=*]
\item If $q$ is nondegenerate and is either definite or of dimension at
most two, then $V_A=V_\Z$.
\item If $q$ is degenerate, every $\mathbf a\in V_\Z$ lies on an
injective affine-linear family $\mathbf a+t\mathbf v$, $t\in\J$,
in $V_A$, with every $t\ne0$ giving an infinite point.
\item If $q$ is nondegenerate, indefinite, and $n\geq3$, every
$\mathbf a\in V_\Z$ lies on an injective quadratic polynomial family
$\mathbf a+t\mathbf v+t^2\mathbf w$, $t\in\J$, in $V_A$, again
with every $t\ne0$ giving an infinite point.
\end{enumerate}
In alternatives \textup{(ii)} and \textup{(iii)}, the constant tuple of
every point in the family is $\mathbf a$.
\end{theorem}

We first prove the geometric construction required for the indefinite
case. All vectors and bilinear forms in its statement are ordinary real
objects; only the parameter is later allowed to be surreal.

\begin{lemma}[A polynomial quadratic isometry]\label{lem:transvection}
Let $B$ be a symmetric bilinear form over $\R$, with $q(x)=B(x,x)$.
Suppose $e,u$ satisfy
\[
 q(e)=0,\qquad B(e,u)=0,\qquad k=q(u)\ne0.
\]
For a scalar parameter $t$ define
\begin{equation}\label{eq:transvection}
 T_t(x)=x+t\bigl(B(x,e)u-B(x,u)e\bigr)
       -\frac{k}{2}t^2B(x,e)e.
\end{equation}
Over every field extension of $\R$,
\[
 q(T_t(x))=q(x),\qquad T_sT_t=T_{s+t},\qquad T_t^{-1}=T_{-t}.
\]
If $\alpha=B(a,e)\ne0$, the curve $t\mapsto T_t(a)$ is injective.
\end{lemma}
\begin{proof}
Put $\alpha=B(x,e)$ and $\beta=B(x,u)$. In
\eqref{eq:transvection}, the added vector is
\[
 t\alpha u-\left(t\beta+\frac{k}{2}t^2\alpha\right)e.
\]
Its scalar product with $x$, multiplied by two, is
$-kt^2\alpha^2$, because the two linear terms cancel. Its own quadratic
value is $kt^2\alpha^2$, since $e$ is isotropic and orthogonal to $u$.
Their sum is zero, proving preservation of $q$.

For the group law, set $N(x)=B(x,e)u-B(x,u)e$. Then
\[
 N^2(x)=-kB(x,e)e,\qquad N^3=0,
 \qquad T_t=\id+tN+\frac{t^2}{2}N^2.
\]
Multiplying these finite polynomials in $N$ gives $T_sT_t=T_{s+t}$,
and hence the inverse formula.

Finally, for $\alpha=B(a,e)\ne0$ and $\beta=B(a,u)$,
\[
 T_t(a)=a+t v+t^2 w,
 \quad v=\alpha u-\beta e,
 \quad w=-\frac{k\alpha}{2}e.
\]
The real vectors $v,w$ are linearly independent: $w$ is a nonzero
multiple of $e$, whereas $v$ has nonzero $u$-component, and $u,e$ are
independent. Their independence persists over a field extension because
some $2\times2$ real minor is nonzero. If $T_s(a)=T_t(a)$ and $s\ne t$,
then $v+(s+t)w=0$, contradicting that independence.
\end{proof}

\begin{lemma}\label{lem:nullchoices}
If $q$ is a nondegenerate indefinite real quadratic form of dimension at
least three and $q(a)\ne0$, there exist $e,u$ as in
\cref{lem:transvection} with $B(a,e)\ne0$.
\end{lemma}
\begin{proof}
The real isotropic vectors span the entire space. To see this, use a
real basis in which $q$ has positive unit vectors $p_i$ and negative
unit vectors $n_j$. Both types exist by indefiniteness, and each
$p_i+n_j$ and $p_i-n_j$ is isotropic. Taking sums and differences
recovers every basis direction. Since the nonzero functional $B(a,-)$
cannot vanish on a spanning set, some isotropic $e$ has $B(a,e)\ne0$.

Nondegeneracy gives $f_0$ with $B(e,f_0)=1$. Replacing it by
$f=f_0-\tfrac12q(f_0)e$ makes $q(f)=0$ while retaining $B(e,f)=1$.
The plane spanned by $e,f$ is nondegenerate. Its orthogonal complement
is nondegenerate and has dimension $n-2\geq1$, so it contains $u$ with
$q(u)\ne0$. This $u$ is orthogonal to $e$, as required.
\end{proof}

\begin{proof}[Proof of \cref{thm:quadclassification}]
The empty-set assertion is \cref{thm:transfer}.

Suppose first that $q$ is definite and $\mathbf x\in V_A$ has an
infinite coordinate. Let $\gamma>0$ be the largest of the leading
exponents of its nonzero coordinates. Let $r_i$ be the coefficient of
$\omega^\gamma$ in $x_i$, using zero when that exponent is absent.
The real vector $\mathbf r$ is nonzero. The coefficient of
$\omega^{2\gamma}$ in $q(\mathbf x)$ is exactly $q(\mathbf r)$,
which is nonzero by definiteness. This contradicts $q(\mathbf x)=c$.
Thus every coordinate is finite and hence an ordinary integer.
In dimension two, nondegeneracy means that the binary quadratic has
two distinct linear factors over $\C$, so \cref{thm:binary} applies.
A nondegenerate one-dimensional form is definite. This proves (i).

For (ii), choose a nonzero real vector $v\in\rad(B)$. For every $t$,
$q(a+tv)=q(a)=c$. If $t\in\J$, the coordinates of $tv$ belong to
$\J$, so $a+tv\in A^n$ with constant tuple $a$. The map is injective
because $v\ne0$. For $t\ne0$, some coordinate has positive degree.

For (iii), choose $e,u$ by \cref{lem:nullchoices} and use
\cref{lem:transvection}. Each coordinate of $T_t(a)$ is an integer
constant plus a real multiple of $t$ plus a real multiple of $t^2$.
Thus it belongs to $A$ and has constant coefficient $a_i$ for
$t\in\J$. Injectivity has already been proved. If $t\ne0$, its
leading degree is positive. In a coordinate where $w_i\ne0$, the
quadratic term has degree $2\ddeg(t)$, strictly larger than the linear
term's degree, and cannot cancel. Therefore the point is infinite.
\end{proof}

\subsection{An explicit infinite hyperboloid family}

For positive ordinary integers $D,E$ and $t\in\J$, set
\begin{equation}\label{eq:hyperboloid}
 x=1+\frac E2t^2,\qquad
 y=\frac{E}{2\sqrt D}t^2,\qquad z=t.
\end{equation}
All three coordinates are omnific integers, and a direct expansion gives
\[
 x^2-Dy^2-Ez^2=1.
\]
Indeed, the two quartic terms cancel, as do the two quadratic terms.
For $t\ne0$ every coordinate is infinite. Nevertheless the tuple is
unimodular, since
\begin{equation}\label{eq:hyperunimodular}
 1=x-\left(\frac E2t\right)z,
 \qquad \frac E2t\in A.
\end{equation}
This illustrates a particularly sharp contrast:
\[
 x^2-Dy^2=1\ \text{has only ordinary points},
 \qquad x^2-Dy^2-Ez^2=1\ \text{has infinite unimodular points}.
\]
There is no appeal to convergence, analytic continuation, or infinite
iteration in this construction; every expression is a finite polynomial.

\subsection{What the classification does and does not say}

The classification concerns the presence of genuinely new points on a
\emph{represented nonzero level}. It does not decide whether the
integer level is represented in the first place. Nor does it classify all
infinite points: a displayed curve is a rigorously constructed subset of
the full fiber.

For zero levels, homogeneity gives a different mechanism. An indefinite
real quadratic form has real isotropic directions, and multiplying any
such direction by an element of $\J$ produces omnific points with zero
constant tuple. Such points can exist even when there is no nonzero
ordinary integer isotropic vector.

\section{Defining the ordinary integers by one quartic equation}\label{sec:guard}

The retraction $\ct$ detects ordinary coefficients but is not a symbol
in the ring language. We now show that the embedded copy of $\Z$ can be
detected by a single explicit polynomial equation with existential
auxiliary variables.

\subsection{Pell solutions as finite bounds}

Define ordinary positive Pell solutions by
\begin{equation}\label{eq:pellsequence}
 u_k+v_k\sqrt2=(3+2\sqrt2)^k,\qquad k\geq0.
\end{equation}
Then $u_k,v_k\in\Z$, $u_k^2-2v_k^2=1$, and $u_k$ tends to infinity
in the ordinary real sense. The inverse of $3+2\sqrt2$ is
$3-2\sqrt2\in(0,1)$, so
\[
 u_k=\frac{(3+2\sqrt2)^k+(3-2\sqrt2)^k}{2}.
\]
The recurrence
\[
 u_0=1,\quad u_1=3,\quad u_{k+2}=6u_{k+1}-u_k
\]
implies
\begin{equation}\label{eq:pellmod8}
 u_{2j}\equiv1\pmod8,\qquad
 u_{2j+1}\equiv3\pmod8.
\end{equation}
Both parity subsequences are unbounded.

We use the classical Legendre--Gauss three-square theorem: a nonnegative
ordinary integer $N$ is a sum of three integer squares if and only if
it is not of the form $4^a(8b+7)$, with $a,b\geq0$. This theorem is
imported, not reproved here; see the mathematical exposition in
\cite{Shiu} and the formal development in \cite{ThreeSquares}.

\begin{lemma}[A three-square gap below a Pell coordinate]\label{lem:pellgap}
For every ordinary nonnegative integer $S$, there are ordinary integers
$u,v,w_1,w_2,w_3$ such that
\[
 u^2-2v^2=1,\qquad u=S+w_1^2+w_2^2+w_3^2.
\]
\end{lemma}
\begin{proof}
Choose the residue of $u$ modulo $8$ according to the following table.
Each choice is realized by an unbounded subsequence in
\eqref{eq:pellsequence}.
\begin{center}
\begin{tabular}{c*{8}{r}}
\toprule
$S\bmod8$&0&1&2&3&4&5&6&7\\
chosen $u\bmod8$&1&3&3&1&1&3&1&1\\
$(u-S)\bmod8$&1&2&1&6&5&6&3&2\\
\bottomrule
\end{tabular}
\end{center}
Take a sufficiently large $u$ in that subsequence so that $N=u-S>0$.
Its residue is one of $1,2,3,5,6$. Hence it is neither divisible by $4$
nor congruent to $7$ modulo $8$, so it cannot be $4^a(8b+7)$.
The three-square theorem supplies the $w_j$. The corresponding Pell
coordinate supplies $v$.
\end{proof}

\subsection{The uniform five-auxiliary-variable guard}

For each ordinary $n\geq1$, define the integer polynomial
\begin{equation}\label{eq:guardpoly}
 \begin{split}
 G_n(\mathbf X,U,V,\mathbf W)
 ={}&(U^2-2V^2-1)^2\\
 &+\left(U-\sum_{i=1}^nX_i^2-
                   W_1^2-W_2^2-W_3^2\right)^2.
 \end{split}
\end{equation}
It has total degree four and $n+5$ variables.

\begin{theorem}[Uniform quartic integrality guard]\label{thm:guard}
For every $\mathbf x\in A^n$,
\begin{equation}\label{eq:guarddef}
 \mathbf x\in\Z^n
 \quad\Longleftrightarrow\quad
 \exists u,v,w_1,w_2,w_3\in A\quad
 G_n(\mathbf x,u,v,w_1,w_2,w_3)=0.
\end{equation}
In fact, every coordinate of every zero of $G_n$ in $A^{n+5}$ is an
ordinary integer. The number of auxiliary variables is five,
independently of $n$.
\end{theorem}
\begin{proof}
Suppose $G_n=0$ in $A$. A sum of two squares in an ordered field can
vanish only when both squares vanish. Thus
\begin{equation}\label{eq:guardrelations}
 u^2-2v^2=1,
 \qquad u=\sum_i x_i^2+w_1^2+w_2^2+w_3^2.
\end{equation}
By Pell rigidity, $u,v\in\Z$. The second equation gives $u\geq0$
and bounds the square of each $x_i$ and $w_j$ by the ordinary number
$u$. Consequently every $x_i,w_j$ is finite. By
\cref{prop:units}, each is an ordinary integer.

Conversely, suppose $\mathbf x\in\Z^n$ and put
$S=\sum_i x_i^2\in\N$. Choose an ordinary Pell solution and three
squares from \cref{lem:pellgap}. They satisfy both equations in
\eqref{eq:guardrelations}, and hence make $G_n$ vanish.
\end{proof}

The theorem is uniform in tuple length, but does not claim that five
auxiliary variables are optimal. A slightly simpler version using four
square variables instead of three follows from Lagrange's four-square
theorem and arbitrary sufficiently large Pell coordinates; the modular
argument above removes one auxiliary variable.

For $x=1$, one witness is $u=3$, $v=2$, and
$(w_1,w_2,w_3)=(1,1,0)$. For $x=2$, a witness is $u=17$, $v=12$,
and $(w_1,w_2,w_3)=(3,2,0)$. The assertions for all integers rely on
\cref{lem:pellgap}, not on these examples or on a finite computation.

\subsection{Definability consequences and a logical distinction}

Let $\Int(x)$ abbreviate the existential formula obtained from $G_1$.
Then $\Int(A)=\Z$. Thus ordinary integer arithmetic is present not
only as the quotient $A/\J$ but as an existentially definable subring of
$A$ itself.

There is no conflict with \cref{prop:notdiophantine}. The copy of $\Z$
is stable under $\ct$, whereas the full positive cone and the nonzero
locus are not. Nor does this make $A$ an elementary extension of $\Z$:
the formula $\Int(x)$ holds for every element of $\Z$ but fails for
$\omega$ in $A$.

\begin{corollary}\label{cor:Jdefinable}
The ideal $\J$ and the graph of $\ct$ are first-order definable in the
pure ring structure of $A$, with no nonstandard parameters. Explicitly,
\begin{align*}
 j\in\J\quad\Longleftrightarrow\quad
 &\forall n\,
  \bigl((\Int(n)\land n\ne0)\ \Longrightarrow\
                 \exists y\ (j=ny)\bigr),\\
 m=\ct(x)\quad\Longleftrightarrow\quad
 &\Int(m)\land(x-m\in\J).
\end{align*}
\end{corollary}
\begin{proof}
The first formula restricts the divisor variable to the ordinary nonzero
integers, then applies \eqref{eq:divkernel}. The second expresses the
unique decomposition $x=m+j$ with $m\in\Z$, $j\in\J$.
\end{proof}

The displayed definition of $\J$ uses a universal quantifier and an
implication. We do not claim that it is an existential or Diophantine
definition of $\J$.

There is also a concrete failure of a familiar arithmetic sentence. The
positive omnific integer $\omega-1$ is not a sum of any fixed finite
number of omnific squares. Applying $\ct$ to such a proposed identity
would give
\[
 -1=m_1^2+\cdots+m_r^2\qquad(m_i\in\Z),
\]
which is impossible. In particular, the ordinary assertion that every
nonnegative integer is a sum of four squares fails for $A$'s full
nonnegative cone. This is compatible with the classical fact that
integer parts of real closed fields satisfy Open Induction
\cite{Shepherdson,LM}; Open Induction is not full ordinary arithmetic.

\subsection{A reusable abstract principle}

The proof of \cref{thm:guard} needs less than all surreal numbers.
Let $D$ be any discretely ordered ring containing $\Z$, and suppose
that every solution in $D$ of $U^2-2V^2=1$ is ordinary. Every element
of $D$ bounded by an ordinary integer is itself ordinary: discreteness
and a finite comparison with the integers in the bounding interval prove
this. The identical proof therefore shows that $G_n$ defines $\Z^n$
in $D^n$. Our contribution is an explicit guard under a transparent
Pell-rigidity hypothesis, together with a proof that the omnific rings
under consideration satisfy that hypothesis.

\section{Infinite families, primitivity, and projective points}\label{sec:primitive}

\subsection{Polynomial curves through integer points}

\begin{proposition}[Finite polynomial lifting]\label{prop:polylift}
Let $F_1,\ldots,F_r\in\R[X_1,\ldots,X_n]$, and let
$P_1(T),\ldots,P_n(T)\in\R[T]$ satisfy
\[
 F_j(P_1(T),\ldots,P_n(T))=0\quad(1\leq j\leq r),
 \qquad P_i(0)\in\Z\quad(1\leq i\leq n).
\]
Then $\mathbf P(t)\in A^n$ solves the system for every $t\in\J$, with
constant tuple $\mathbf P(0)$. If some $P_i$ is nonconstant, every
$t\in\J\setminus\{0\}$ produces an infinite point.
\end{proposition}
\begin{proof}
Every nonconstant term of $P_i(t)$ lies in $\J$, so the constant
coefficient is exactly the integer $P_i(0)$. Polynomial identities hold
under substitution into any field extension. If $P_i$ has degree $d>0$
and $\ddeg(t)=\gamma>0$, its top term has degree $d\gamma$, strictly
larger than every other term's degree, so cannot cancel.
\end{proof}

The finiteness of the polynomials matters. Evaluating an infinite real
power series at an infinite omnific argument is not licensed by this
proposition. For instance, the Taylor expansion of $\sqrt{1+T}$ at
zero is not an admissible way to assert that $\sqrt{1+\omega}$ lies
in $A$.

\subsection{Infinite solutions of homogeneous equations}

\begin{theorem}[Homogeneous cone criterion]\label{thm:cone}
Let $F_1,\ldots,F_r\in\R[X_1,\ldots,X_n]$ be homogeneous polynomials
of positive degrees. Their common zero locus has a nonzero point in
$A^n$ if and only if it has a nonzero point in $\R^n$.
\end{theorem}
\begin{proof}
A real nonzero common zero $\mathbf a$ gives the omnific common zero
$\omega\mathbf a\in\J^n$ by homogeneity.

Conversely, take a nonzero common zero $\mathbf x\in A^n$. If every
coordinate is constant, it is already a real point. Otherwise let
$\gamma>0$ be the maximum of the coordinates' leading degrees and let
$r_i=[\omega^\gamma]x_i$. The vector $\mathbf r\in\R^n$ is nonzero.
For a homogeneous $F_j$ of degree $d_j$, the coefficient at
$\omega^{d_j\gamma}$ in $F_j(\mathbf x)$ is $F_j(\mathbf r)$: a
product can attain that exponent only by taking the exponent $\gamma$
in every factor. Since $F_j(\mathbf x)=0$, this coefficient vanishes.
Thus $\mathbf r$ is a nonzero real common zero.
\end{proof}

For example, for every ordinary integer $n\geq2$,
\begin{equation}\label{eq:fermatscaling}
 x=t,
 \qquad y=t,
 \qquad z=2^{1/n}t,
 \qquad 0<t\in\J
\end{equation}
gives a positive infinite solution of $x^n+y^n=z^n$. The irrational
real coefficient of $t$ is permitted by the omnific definition.
All three constant coefficients are zero, so these families carry no
nonzero ordinary integer solution under the retraction.

All coordinates in \eqref{eq:fermatscaling} are divisible by $2$ in
$A$, since they belong to $\J$. By \cref{thm:commondivisor}, they
also have a common infinite monomial divisor. Hence these are
not primitive solutions. One must not cite them as resolving a primitive
Fermat problem.

\subsection{Primitive versus unimodular tuples}

\begin{definition}\label{def:primitive}
A nonzero tuple over $A$ is \emph{primitive} if it has no common
nonunit divisor. It is \emph{unimodular} if its coordinates generate the
unit ideal, equivalently if an $A$-linear combination of them is $1$.
\end{definition}

Unimodularity implies primitivity. The converse is not being assumed in
this non-PID ring. We use explicit identities whenever a tuple is claimed
to be unimodular.

\begin{proposition}[Infinite unimodular Pythagorean triples]\label{prop:pythagoras}
For every $t\in\J$, put
\[
 x=t^2-1,\qquad y=2t,\qquad z=t^2+1.
\]
Then $x^2+y^2=z^2$ and the tuple is unimodular. For $t\ne0$ it is a
genuinely infinite triple.
\end{proposition}
\begin{proof}
The identity is the ordinary polynomial identity
$(t^2-1)^2+4t^2=(t^2+1)^2$. Put $a=t^2/2\in\J\subseteq A$.
Then
\begin{equation}\label{eq:pythbezout}
 a z-(a+1)x=\frac{t^2}{2}(z-x)-x=t^2-(t^2-1)=1.
\end{equation}
Thus even the pair $(x,z)$ generates the unit ideal. Positive leading
degree gives the assertion about infinite size.
\end{proof}

The constant tuple is $(-1,0,1)$, not the zero tuple. That is exactly
where the common-divisor obstruction for $\J^3$ ceases to apply.

\subsection{Which real projective directions have primitive representatives?}

A point $[r_0:\cdots:r_n]\in\mathbb P^n(\R)$ is \emph{rational} if
it has homogeneous coordinates in $\Q$, equivalently a nonzero integer
coordinate tuple.

\begin{theorem}[Rationality of primitive constant directions]\label{thm:realray}
For a real projective point $p=[r_0:\cdots:r_n]$, the following are
equivalent:
\begin{enumerate}[label=\textup{(\roman*)},leftmargin=*]
\item $p$ is rational;
\item $p$ has a unimodular representative in $A^{n+1}$;
\item $p$ has a primitive representative in $A^{n+1}$.
\end{enumerate}
\end{theorem}
\begin{proof}
For (i)$\Rightarrow$(ii), choose an integer representative whose ordinary
integer gcd is one. The usual integer B\'ezout identity makes the tuple
unimodular over $A$ as well. Implication (ii)$\Rightarrow$(iii) follows
from the definition.

For (iii)$\Rightarrow$(i), write a primitive representative as
$z_i=\lambda r_i$ with $\lambda\in\No^\times$. Choose $r_j\ne0$.
Since $z_j\in A\subseteq\BR$ and division by the nonzero real scalar
$r_j$ preserves $\BR$, we have $\lambda\in\BR$.
Put $c=\ct(\lambda)\in\R$. Then
\[
 \ct(z_i)=c r_i\in\Z\qquad(0\leq i\leq n).
\]
If $c=0$, every $z_i$ belongs to $\J$, so $2$ divides every
coordinate in $A$, contradicting primitivity. Hence $c\ne0$, and $(cr_0,\ldots,cr_n)$ is a nonzero
integer representative of $p$.
\end{proof}

For example, $[\sqrt2:1]$ has omnific representatives such as
$(\sqrt2\,\omega,\omega)$, but has no primitive omnific
representative. In contrast, each rational real direction has an ordinary
coprime integer representative. The theorem concerns real projective
points; it is not a classification of all points of $\mathbb P^n(\No)$.

\section{Obstructions from leading terms and bounded geometry}\label{sec:obstructions}

\subsection{Initial equations are necessary, not sufficient}

Let $F(\mathbf X)=\sum_{\alpha}a_\alpha\mathbf X^\alpha$ be a
nonzero polynomial over $A$, and suppose all coordinates of a proposed
solution $\mathbf x\in A^n$ are nonzero. For nonzero coefficients put
\[
 \lambda_\alpha=\ddeg(a_\alpha),\quad
 c_\alpha=\lc(a_\alpha),\quad
 \gamma_i=\ddeg(x_i),\quad r_i=\lc(x_i).
\]
The largest possible exponent in $F(\mathbf x)$ is
\[
 M=\max_\alpha\left(\lambda_\alpha+
                                \sum_i\alpha_i\gamma_i\right).
\]
If $F(\mathbf x)=0$, at least two monomials must attain $M$, and
\begin{equation}\label{eq:initial}
 \sum_{\lambda_\alpha+\sum_i\alpha_i\gamma_i=M}
 c_\alpha\prod_i r_i^{\alpha_i}=0.
\end{equation}
This follows by extracting the coefficient at $M$. Variables assigned
zero can first be substituted out of the polynomial.

For standard coefficients this is the usual leading-weight obstruction,
now with surreal weights. It must be used together with the
constant-coefficient equations. Neither set of necessary conditions is
sufficient in general.

For example, $Y^2=\omega^2+1$ passes the constant test
$\ct(Y)^2=1$ and the leading test $\gamma_Y=1$, $r_Y=\pm1$.
Nevertheless, there is no solution in $A$, by the unit argument in
\cref{prop:notnormal}. The missing information resides in the full tail
of a potential root.

\subsection{No real point at infinity implies rigidity}

\begin{theorem}[Leading-homogeneous obstruction]\label{thm:infinity}
Let $F_1,\ldots,F_r\in\Z[X_1,\ldots,X_n]$ be nonconstant polynomials,
and let $H_j$ be the top homogeneous part of $F_j$. If the $H_j$ have
no common nonzero real zero, then every common zero of the $F_j$ in
$A^n$ lies in $\Z^n$.
\end{theorem}
\begin{proof}
Suppose a common zero has an infinite coordinate. Choose its maximum
coordinate degree $\gamma>0$ and nonzero leading vector
$\mathbf r$ as in \cref{thm:cone}. If $d_j=\deg F_j$, the coefficient
at $d_j\gamma$ in $F_j(\mathbf x)$ is $H_j(\mathbf r)$.
Every term of lower ordinary total degree has exponent strictly less
than $d_j\gamma$, and therefore cannot affect that coefficient.
Thus every $H_j(\mathbf r)$ is zero, contrary to the hypothesis.
All coordinates are consequently finite and hence ordinary integers.
\end{proof}

The hypothesis can be viewed as a sufficient absence-of-real-directions-
at-infinity condition for the chosen system. Its converse is false.
A Pell equation has real directions at infinity when $D>0$, but its
nonzero fibers are rigid by factorization. Initial-form obstructions and
constant-product obstructions detect different mechanisms.

\subsection{Positive bounds and orthogonal matrices}

\begin{proposition}\label{prop:squaresbound}
If $a_1,\ldots,a_n$ are positive ordinary real numbers,
$x_1,\ldots,x_n\in A$, and
\[
 \sum_i a_i x_i^2=C\in\R,
\]
then every $x_i$ is an ordinary integer. In particular, an omnific point
on a sphere of ordinary finite radius is an ordinary integer point.
\end{proposition}
\begin{proof}
Each $a_i x_i^2$ is nonnegative and at most $C$. If $C<0$ there is no
solution; otherwise this gives an ordinary real bound on each $|x_i|$.
Use \cref{prop:units}.
\end{proof}

\begin{corollary}\label{cor:orthogonal}
The orthogonal matrices over $A$ are exactly the ordinary signed
permutation matrices. Thus $\OO_n(A)$ is the finite group of order
$2^n n!$, and the determinant-one subgroup consists of its signed
permutation matrices of determinant one.
\end{corollary}
\begin{proof}
For an orthogonal matrix, the sum of the squares in each row is $1$.
Every entry is therefore ordinary. An integer row of squared norm one
has exactly one nonzero entry, equal to $1$ or $-1$. Orthogonality of
different rows forces those entries into different columns. Conversely,
each signed permutation matrix is orthogonal.
\end{proof}

By contrast, $\GL_2(A)$ contains all the infinite unipotent matrices
\[
 \begin{pmatrix}1&t\\0&1\end{pmatrix},\qquad t\in A.
\]
The Euclidean orthogonal group is rigid because of positivity, not because
all matrix groups over $A$ are finite or constant.

\begin{corollary}[Symmetric matrix rigidity]\label{cor:symmetric}
Let $M$ be a symmetric matrix over $A$. If $\tr(M^2)$ is an ordinary
real constant, then every entry of $M$ is an ordinary integer. In
particular, this holds if its characteristic polynomial has ordinary
integer coefficients.
\end{corollary}
\begin{proof}
For symmetric $M$,
\[
 \tr(M^2)=\sum_i m_{ii}^2+2\sum_{i<j}m_{ij}^2,
\]
so \cref{prop:squaresbound} applies. The first two characteristic
coefficients give
$\tr(M^2)=(\tr M)^2-2e_2(M)$ when the size is at least two;
for size one the assertion is immediate. Hence a constant characteristic
polynomial makes this trace constant.
\end{proof}

\subsection{Representative equations}

\begin{center}
\renewcommand{\arraystretch}{1.22}
\begin{tabular}{p{0.34\textwidth}p{0.56\textwidth}}
\toprule
Equation or constraint & Exact conclusion over $A$\\
\midrule
$2x+3y=1$ & $x=-1+3t$, $y=1-2t$, $t\in A$.\\
$x^2+y^2=N\in\Z$ & Only ordinary integer solutions.\\
$x^2-Dy^2=c$, $Dc\ne0$ & Only ordinary integer solutions.\\
$x^2-Dy^2=0$, $D>0$ nonsquare & $(x,y)=(\pm\sqrt D\,t,t)$, $t\in\J$.\\
$x^4-y^6=c\ne0$ & Only ordinary integer solutions.\\
$x^2-y^2-z^2=1$ & Infinite unimodular family from \eqref{eq:hyperboloid}.\\
$x^2+y^2=z^2$ & Infinite unimodular family from \cref{prop:pythagoras}.\\
$x^n+y^n=z^n$, $n\geq2$ & Positive infinite solutions in \eqref{eq:fermatscaling}; these are not primitive.\\
$y^2=\omega^2+1$ & No solution, despite consistent constant and leading tests.\\
$Q^\mathsf{T}Q=I$ & $Q$ is an ordinary signed permutation matrix.\\
\bottomrule
\end{tabular}
\end{center}

\section{Fixed Hahn workspaces and possible formalization}\label{sec:formal}

\subsection{Which results are local in the exponent group?}

For a set-sized ordered abelian group $\Gamma$, work in a Hahn field
whose monomials are written $\omega^\gamma$ with reverse-well-ordered
supports, and define
\[
 A_\Gamma=\Z+
 \left\{\sum_{\gamma>0}r_\gamma\omega^\gamma\right\},
 \qquad
 B_\Gamma=\R+
 \left\{\sum_{\gamma>0}r_\gamma\omega^\gamma\right\}.
\]
The displayed sums mean elements of the ambient Hahn field with the
specified support restrictions. This notation is a convention, not an
identification of every abstract group element with a surreal exponent.
When an appropriate ordered group embedding into $\No$ has been fixed,
it gives the corresponding surreal normal-form workspace.

The retraction, finite quotients, equational transfer, Smith criterion,
constant-product rigidity, norm rigidity, quadratic constructions, and
quartic guard all apply to $A_\Gamma$. Their proofs use finite algebra,
leading exponents, and real coefficients, not the availability of a new
exponent outside $\Gamma$. If $\Gamma\ne0$, its positive cone provides
nonconstant parameters for the infinite-family statements. The case
$\Gamma=0$ is simply $A_\Gamma=\Z$.

By contrast, the common-divisor theorem for an arbitrary set of purely
infinite elements and the monomial clearing theorem for an arbitrary set
of full-field elements require bounds that may not exist in $\Gamma$.
For instance, in $\Gamma=\R$, the set $\{1/k:k\geq1\}$ has no
strictly positive real lower bound. In the full surreal group it does.

One must therefore distinguish a proof that stays inside a fixed workspace
from a proof that enlarges it. In particular, the equality
$\Frac(A)=\No$ and the no-set-of-generators statement for $\J$ are
full-class conclusions. They are not asserted for every $A_\Gamma$.

\subsection{A dependency-oriented formalization plan}

The repository already has normal-form and Hahn algebra infrastructure,
coefficient extraction, and a carefully stated boundary between concrete
surreal results and workspace results \cite{Repo}. The following is a
\emph{proposed} extension plan. The names below are suggested module
names, not claims that corresponding files or declarations already exist.

\begin{center}
\renewcommand{\arraystretch}{1.2}
\begin{tabular}{p{0.34\textwidth}p{0.55\textwidth}}
\toprule
Proposed module & Main obligation and dependency\\
\midrule
\texttt{Omnific.Basic} & Define the nonnegative-exponent subring and the integer-constant subring; prove the retraction and finite-element criterion.\\
\texttt{Omnific.Residue} & Prove $nA=\J+n\Z$, finite quotients, and preservation of $\J$.\\
\texttt{Omnific.Rigidity} & Prove units of the real and complex support rings are constants; derive product, binary-form, and norm rigidity.\\
\texttt{Omnific.Quadratic} & Verify the finite bilinear identity for $T_t$, the isotropic-vector choices, and the level classification.\\
\texttt{Omnific.IntegralityGuard} & Combine Pell rigidity, ordinary Pell residues, and an explicitly imported three-square theorem.\\
\texttt{Omnific.Global} & Use the genuine small-cut theorem for support bounds, common monomial divisibility, and global denominator clearing.\\
\bottomrule
\end{tabular}
\end{center}

A useful separation is to prove the finite algebraic identities first over
arbitrary suitable fields, then instantiate them in a normal-form model.
The transvection identity, for example, is finite bilinear algebra; it has
no dependence on surreal birthday recursion. The only difficult
number-theoretic input to the guard is the classical three-square theorem,
not an unproved claim about omnific factorization.

No Lean files are included here because this manuscript has not been
formalized or checked against a built repository. A proposed module map
must not be confused with a completed formal verification.

\section{Further questions and conclusions}\label{sec:questions}

\subsection{Questions not settled by this article}

The following are precise continuation problems. They are not all claimed
to be open in the literature; the point is to isolate what the present
proofs do not decide.

\paragraph{Primitive higher-degree Fermat fibers.}
The scaled solutions in \eqref{eq:fermatscaling} are never primitive.
For a fixed $n\geq3$, classify primitive or unimodular omnific solutions
of $X^n+Y^n=Z^n$ with all coordinates nonzero. The binary nonzero-fiber
theorem does not directly answer a three-variable homogeneous question.

\paragraph{Coprime exponent equations and affine curves.}
Determine which curves such as $Y^2=X^3+1$ have genuinely infinite
omnific points. The leading equation gives only a necessary condition;
it does not decide whether an algebraic expansion has an inadmissible
negative tail. A satisfactory criterion should separate polynomial
parametrization, torus rigidity, and genuinely algebraic branches.

\paragraph{The arithmetic of the pure-infinite ideal.}
We have given a first-order definition of $\J$, but not an existential
one. Is there an existential definition over $\Z$? More generally,
which fibers or subideals of $\J$ are Diophantine? Retraction stability
is necessary for such definitions but far from sufficient.

\paragraph{Complexity of the integrality guard.}
Can one reduce the number of auxiliary variables below five, retain a
single quartic equation, or prove a meaningful lower bound? The present
construction establishes an explicit upper bound only. The role of the
three-square theorem suggests testing other rigid arithmetic families as
finite bounds.

\paragraph{Infinite coefficients and effective fragments.}
For polynomials with finitely represented omnific coefficients, what can
be decided about roots in a chosen support class? The equation
$Y^2=\omega^2+1$ shows why a solver needs more than constant and leading
conditions. Any effective claim must specify coefficients, exponent
representations, and whether workspace enlargement is permitted.

\subsection{Synthesis}

The arithmetic of $A$ is governed by three simultaneous facts:
\[
 A/\J=\Z,\qquad A^\times=\{\pm1\},\qquad\Frac(A)=\No.
\]
The first makes ordinary congruences and equation-existence questions
look classical. The second imposes powerful rigidity on nonzero product
fibers. The third makes projective representation extraordinarily flexible,
while saying little about primitive representatives.

The difference between a Pell conic and a three-variable indefinite
quadric is therefore structural, not a matter of choosing a sufficiently
large number. Binary factorization forces constants; an additional
indefinite dimension supports finite polynomial isometries. The quartic
guard uses the rigid side of this dichotomy to recover ordinary integers
from the much larger omnific ring. These mechanisms provide a workable
starting theory of omnific--Diophantine geometry without assuming unique
factorization, a Euclidean termination principle, or full transfer from
ordinary arithmetic.

\appendix
\section{Proof dependencies and verification boundaries}\label{app:dependencies}

The foundation imported from surreal theory is the existence, uniqueness,
order, and multiplication of set-supported Conway normal forms, the
set-cut property, and real closedness of $\No$. All other main deductions
are proved from these and ordinary finite algebra, except where an
external number-theoretic theorem is explicitly cited.

The principal dependency chains are:
\begin{align*}
 &\text{normal forms}\ \Longrightarrow\ \ct\text{ is a retraction}
 \ \Longrightarrow\ \text{equational transfer and finite quotients};\\
 &\text{normal forms}\ \Longrightarrow\ \text{constant units in }\BC
 \ \Longrightarrow\ \text{binary and norm rigidity};\\
 &\text{Pell rigidity} + \text{ordinary three-square theorem}
 \ \Longrightarrow\ \text{quartic integrality guard};\\
 &\text{real quadratic linear algebra} + \text{finite support-ring algebra}
 \ \Longrightarrow\ \text{quadratic-level dichotomy};\\
 &\text{normal forms} + \text{set-cut bounds}
 \ \Longrightarrow\ \text{global common divisors and denominator clearing}.
\end{align*}

The included \texttt{verification.py} checks the cubic norm identity,
Pythagorean and hyperboloid identities, B\'ezout certificates, finite
matrix instances of the quadratic isometry identity, Pell residue
patterns, and finitely many three-square guard witnesses. It is a
supplement, not a proof of statements involving all normal forms or all
ordinary integers. It does not implement the full class of surreal numbers.

The paper's claims have been reviewed by algebraic derivation and the
finite checks just listed. No independent referee report or proof-assistant
verification is claimed.

\section{Repository and source audit}\label{app:audit}

\subsection{Repository snapshot}

The inspected repository was \repo, pinned to
\begin{center}\small\pin\end{center}
The GitHub connector was used to retrieve the repository tree, root README,
documentation index, trigonometry directory listing, and
\texttt{21-period-arithmetic-source\_audit.md}. The inspected documentation
already discusses normal forms, integer parts, coefficient extraction,
Hahn workspaces, and finite residue arithmetic. These are credited as
background rather than presented as discoveries of this article.

This was a focused source inspection, not a repository checkout, build,
or exhaustive review of every manuscript and Lean declaration. A missing
search match in a large returned tree is not treated as proof that the
repository contains no related result. The article's mathematical proofs
stand on their stated hypotheses, not on an inference from an incomplete
repository search.

\subsection{Primary mathematical sources}

The 2024 L'Innocente--Mantova article was inspected through its arXiv text,
with its opening PDF page also checked as an image. Its normal-form
conventions and distinction between ordinary factorization and generalized
series factorization were especially relevant. Ehrlich--Kaplan's initial
integer-part lemma and Je\v r\'abek's residue-kernel discussion were
checked in the available article text.

For the classical three-square theorem, the primary mathematical paper
\cite{Shiu} and the Archive of Formal Proofs entry \cite{ThreeSquares}
were consulted. For undecidability, the author-hosted text accompanying
Matiyasevich's monograph was consulted \cite{Matiyasevich}. The standard
books of Conway and Gonshor are background references; they were not
newly read in full for this task.

\subsection{A dated factorization announcement}

A September 18, 2026 announcement reports a Lean proof of Conway's
refinement conjecture \cite{RefinementAnnouncement}. Its author explicitly
qualifies the independent-verification status. We did not audit that code
or use the announced result. Accordingly, this article does not repeat
older literature's description of the refinement conjecture as a verified
\emph{currently open} problem, and does not claim that this manuscript
settles it. None of our proofs depends on its present status.

\subsection{Novelty and attribution}

The strongest proposed contributions of this manuscript are the explicit
quadratic-level classification, the uniform five-auxiliary-variable
quartic guard, and the organization of rigidity and primitive-point
obstructions into an omnific--Diophantine framework. Familiar mechanisms
such as coefficient retractions, unit arguments, polynomial
parametrizations, and nilpotent quadratic isometries are not claimed as
new inventions. A focused search and source review do not establish
historical priority; that would require broader expert review.

\begin{thebibliography}{99}
\small\raggedright

\bibitem{Conway}
John H. Conway.
\emph{On Numbers and Games}.
Academic Press, 1976; second edition, A K Peters, 2001.

\bibitem{Gonshor}
Harry Gonshor.
\emph{An Introduction to the Theory of Surreal Numbers}.
London Mathematical Society Lecture Note Series 110.
Cambridge University Press, 1986.

\bibitem{LM}
Sonia L'Innocente and Vincenzo Mantova.
A factorisation theory for generalised power series and omnific integers.
\emph{Advances in Mathematics} \textbf{442} (2024), 109513.
\href{https://doi.org/10.1016/j.aim.2024.109513}{doi:10.1016/j.aim.2024.109513}.
\href{https://arxiv.org/abs/1710.07304}{arXiv:1710.07304}, version 5,
January 22, 2024.

\bibitem{EK}
Philip Ehrlich and Elliot Kaplan.
Surreal ordered exponential fields.
\emph{Journal of Symbolic Logic} \textbf{86}(3) (2021), 1066--1115.
\href{https://doi.org/10.1017/jsl.2021.59}{doi:10.1017/jsl.2021.59}.
\href{https://arxiv.org/abs/2002.07739}{arXiv:2002.07739}.

\bibitem{Jerabek}
Emil Je\v r\'abek.
Rigid models of Presburger arithmetic.
\emph{Mathematical Logic Quarterly} \textbf{65}(1) (2019), 108--115.
\href{https://doi.org/10.1002/malq.201800019}{doi:10.1002/malq.201800019}.
\href{https://arxiv.org/abs/1803.05797}{arXiv:1803.05797}.

\bibitem{Shepherdson}
John C. Shepherdson.
A non-standard model for a free variable fragment of number theory.
\emph{Bulletin de l'Acad\'emie Polonaise des Sciences}
\textbf{12} (1964), 79--86.

\bibitem{Matiyasevich}
Yuri V. Matiyasevich.
\emph{Hilbert's Tenth Problem}.
MIT Press, 1993.
Author-hosted supplementary text, especially Section 5.7:
\url{https://logic.pdmi.ras.ru/~yumat/H10Pbook/par_5_7.htm}.

\bibitem{Shiu}
Peter Shiu.
The three-square theorem of Gauss and Legendre.
\emph{The Mathematical Gazette} \textbf{104}(560) (2020), 209--214.
\href{https://doi.org/10.1017/mag.2020.42}{doi:10.1017/mag.2020.42}.

\bibitem{ThreeSquares}
Anna Danilkin and Lo\"ic Chevalier.
Three Squares Theorem.
\emph{Archive of Formal Proofs}, May 3, 2023.
\url{https://isa-afp.org/entries/Three_Squares.html}.

\bibitem{Repo}
Vladimir Reshetnikov and contributors.
\emph{Surreal}: repository and research documentation.
\url{https://github.com/VladimirReshetnikov/Surreal}.
Snapshot \pin.
See root \texttt{README.md}, \texttt{docs/README.md}, and the
trigonometry source-audit file identified in Appendix~\ref{app:audit}.

\bibitem{RefinementAnnouncement}
\emph{How I Vibed a Proof of Conway's Conjecture}.
Author's announcement on \emph{overreacted}, September 18, 2026.
\url{https://overreacted.io/how-i-vibed-a-proof-of-conways-conjecture/}.
Cited only for the existence and qualified status of the announcement;
not used as a mathematical premise.

\end{thebibliography}
\end{document}
